\chapter*{Prefacio}
\addcontentsline{toc}{chapter}{Prefacio}
\markboth{Prefacio}{Prefacio}

Cuba es uno de los pocos países sobre los que casi nadie escribe con cuidado.
La literatura se divide en dos bibliotecas que no se leen entre sí. Una registra
una isla pequeña y pobre que abolió el analfabetismo, construyó un sistema de
salud muy por encima de su ingreso, envió médicos a cuarenta países y lo hizo
bajo seis décadas de asedio económico por parte de la mayor economía del
planeta. La otra registra un Estado de partido único que encarceló a sus poetas,
criminalizó sus mercados, empujó al exilio a una quinta parte de su gente y
terminó el período con una red eléctrica incapaz de mantener las luces
encendidas. Las dos bibliotecas son exactas. Ninguna de las dos es honesta,
porque cada una está completa solo por omisión.

Este libro intenta sostener ambas a la vez, y lo hace por una razón práctica
antes que académica. La última parte del libro es un diseño: un sistema
sociopolítico y económico propuesto para Cuba a partir de 2026. Un diseño vale
lo que valga su descripción de las condiciones iniciales. Si uno cree en la
primera biblioteca, diseñará como si lo único que anda mal fuera el embargo, y
producirá un plan que fracasa en cuanto se encuentra con un cubano que quiere
abrir una tienda. Si cree en la segunda, diseñará como si lo único que anda mal
fuera el socialismo, y producirá un plan que desmantela un sistema público de
salud que no puede reemplazar, entrega los activos a quien esté más cerca y
reproduce en diez años la desigualdad que hizo posible 1959. Ambos fracasos ya
se han ensayado en otras partes. Ninguno necesita una segunda prueba.

Así que las Partes I y II no son un mero preámbulo. Son el conjunto de
restricciones. La Parte I cubre de 1959 a 1990, los años en que Cuba construyó
su Estado social y lo pagó con un subsidio soviético y una política cerrada. La
Parte II cubre de 1990 a 2026, el largo ajuste a la pérdida de ese subsidio: dos
ciclos completos de liberalizar-bajo-presión y recentralizar-cuando-la-presión-
cede, que terminan en la peor crisis material que la revolución ha enfrentado
---peor en algunas dimensiones que el Período Especial--- y en la mayor
emigración de la historia cubana. La Parte III pregunta qué puede construirse
con lo que queda.

Lo que queda es más de lo que la ruina aparenta. Cuba tiene una población que
sabe leer, calcular y formarse rápido; un establecimiento médico y
biotecnológico con productos reales y una reputación que todavía vale dinero;
algunas de las mejores playas, arrecifes y ciudades coloniales del hemisferio,
subexplotadas; mil cien kilovatios-hora de luz solar por metro cuadrado al año
cayendo sobre un país que hoy quema fueloil importado; una diáspora de unos dos
millones de personas, muchas con recursos, casi todas a un vuelo de distancia;
y, cosa inusual en un país pobre, poca violencia criminal y un orden público que
funciona. Esos son los activos. El argumento del libro es que se
\emph{componen}: que educación más conectividad más un régimen de residencia
para trabajadores del conocimiento móviles constituye una segunda vía
exportadora junto al turismo; que el sol más el almacenamiento convierte el
mayor costo recurrente en divisas del país en un costo doméstico; y que la
biotecnología es el único sector donde Cuba ya posee algo que el mundo pagaría y
nunca se ha organizado para venderlo.

El argumento del libro es también que nada de eso funciona sin la parte
aburrida. Cada uno de los activos enumerados ha estado disponible durante veinte
años y no se ha convertido, y la razón no es solo el embargo. Es que a Cuba le
falta la maquinaria institucional que convierte un activo en una inversión:
contratos exigibles, una moneda que signifique algo, una aduana que no cobre su
tajada, un tribunal ante el cual un extranjero pueda perder limpiamente, un
sistema de compras públicas que publique lo que compró. Un país puede tener sol,
médicos y playas y seguir siendo pobre, y Cuba es la prueba. Por eso los
capítulos más largos de la Parte III no son los de los paneles solares. Son los
del dinero, la propiedad y la corrupción, que es el modo de fallo del que
realmente han muerto las transiciones de este tipo.

Un compromiso de la Parte III no está a discusión en este libro, y conviene
enunciarlo al frente en vez de llegar a él a mitad de camino. El diseño conserva
un Estado de bienestar universal al nivel de Europa occidental ---salud,
educación, pensiones, discapacidad, sostén de ingresos--- y todo lo demás en la
Parte III se ordena alrededor de poder pagarlo. Esto no es sentimentalismo sobre
la revolución. Es la lectura del registro de las Partes I y II: el Estado social
es lo único que esos sesenta y siete años produjeron que sin ambigüedad vale la
pena conservar, es lo que protege a los cubanos más expuestos si una transición
sale mal, y es además un insumo económico y no solo una transferencia, porque el
sistema educativo es precisamente aquello con lo que están hechos la segunda vía
y el sector biotecnológico.

Lo que cambia no es la prestación sino quién la paga. El Estado de bienestar
cubano nunca lo han financiado los cubanos. Lo financiaron los términos de
intercambio soviéticos hasta 1990, el petróleo venezolano en los años dos mil, y
en los intervalos la emisión monetaria y el no mantener nada, que es por lo que
ha estado dos veces al borde del colapso en una sola vida. Una prestación
financiada por una renta externa dura exactamente lo que dura la renta. Los
sistemas europeos que este libro toma como norma se financian gravando una base
doméstica, y por eso sobreviven a los cambios de gobierno, de precio de las
materias primas y de patrón. La propuesta es entonces una estructura europea de
prestaciones sobre una estructura europea de impuestos ---un sistema fiscal de
base amplia y administración sencilla, de un tipo que Cuba nunca ha tenido--- y
el costo de ese compromiso, una carga tributaria muy por encima de la norma
latinoamericana, se declara en los capítulos en vez de esconderse en ellos. Un
país que quiere resultados nórdicos tiene que recaudar ingresos nórdicos, y no
hay versión de esto en que la aritmética resulte agradable.

Escribo como cubano de la diáspora y como físico de formación, lo que marca el
libro de dos maneras que conviene declarar. Es la razón de que la Parte III esté
escrita como un conjunto de reglas institucionales numeradas y no como un
ensayo: un diseño que no puede enunciarse como reglas no es un diseño, y con las
reglas se puede discutir una por una. Es también la razón de que el libro no sea
neutral respecto de los resultados aun cuando intente ser neutral respecto de
los hechos. Quiere que Cuba sea próspera, abierta y dueña de sí. No supone que
esas tres cosas sean automáticamente compatibles, y donde entran en conflicto el
texto lo dice en vez de elegir en silencio.

Nada de lo que sigue es un pronóstico. La perspectiva del título es la del
sentido antiguo de la palabra: una vista de lo que hay por delante desde donde
uno está parado, trazada a escala.
